\documentclass[a4paper]{article}
% Shared preamble for the English probability and statistics documents.
% Based on the original course preamble. Compile with pdfLaTeX.
\usepackage[top=1.5cm,bottom=2cm,left=2cm,right=2cm]{geometry}
\usepackage[T1]{fontenc}
\usepackage{lmodern}
\usepackage[utf8]{inputenc}
\usepackage[english]{babel}

\usepackage{amsmath}
\usepackage{amssymb}
\usepackage{amsthm}
\usepackage{mathrsfs}
\usepackage{mathtools}
\usepackage{etoolbox}
\usepackage{tikz}
\usetikzlibrary{shapes.multipart,arrows.meta}
\usepackage{xcolor}
\usepackage{enumerate}
\usepackage[shortlabels]{enumitem}
\usepackage{hyperref}
\usepackage{version}
\usepackage{bbold}
\usepackage{fontawesome5}
\usepackage{fancyhdr}
\lfoot{}
\usepackage{natbib}
\usepackage{graphicx}
\usepackage{pdfpages}
\usepackage{authblk}
    
\usepackage{blindtext}
\usepackage{multicol}
\usepackage{collcell}
\usepackage{wrapfig}
\usepackage{blkarray}
\usepackage{cancel}
\usepackage{appendix}
\usepackage{float}
\makeatletter
\input{tex/assets/slashbox.sty}
\makeatother
\hypersetup{unicode,colorlinks,citecolor=red,filecolor=black,
  linkcolor=red,urlcolor=cyan}
\setlength{\columnsep}{1cm}
\pagenumbering{gobble}

\def\numtext#1{%
  \ifcase#1\or one\or two\or three\or four\or five\or six\or seven\or eight\or nine\or ten\or
  eleven\or twelve\or thirteen\or fourteen\or fifteen\or sixteen\or seventeen\or eighteen\or nineteen\or twenty\or Lots\fi}

\newcommand{\Byte}[1]{%
  \def\nodecontents{}%
  \foreach[count=\l] \k in {#1} {\xappto{\nodecontents}{\noexpand\nodepart{\numtext{\l}}\k}}%
  \node[rectangle split,rectangle split parts=8,rectangle split horizontal,draw] at (2,2){\nodecontents};}

\newcommand{\D}{{\mathrm{d}}}
\newcommand{\cE}{{\mathcal{E}}}
\newcommand{\cG}{{\mathcal{G}}}
\newcommand{\cX}{{\mathcal{X}}}
\newcommand{\cY}{{\mathcal{Y}}}
\newcommand{\cB}{{\mathcal{B}}}
\newcommand{\cF}{{\mathcal{F}}}
\newcommand{\cO}{{\mathcal{O}}}
\newcommand{\cS}{{\mathcal{S}}}
\newcommand{\cT}{{\mathcal{T}}}
\newcommand{\given}{\, \big| \,}
\newcommand{\cL}{{\mathscr{L}}}
\newcommand{\cD}{{\mathcal{D}}}
\newcommand{\cA}{{\mathscr{A}}}
\newcommand{\cP}{{\mathscr{P}}}
\newcommand{\cR}{{\mathscr{R}}}
\newcommand{\A}{{\mathcal A}}
\newcommand{\one}{\mathbb{1}}
\newcommand{\mathbbm}[1]{\mathbb{#1}}

\providecommand{\C}{}
\renewcommand{\C}{\mathbb C}
\newcommand{\R}{\mathbb R}
\newcommand{\Q}{\mathbb Q}
\newcommand{\N}{\mathbb N}
\newcommand{\T}{\mathbb T}
\newcommand{\F}{\mathbb F}
\renewcommand{\P}{\mathbb P}
\newcommand{\Z}{\mathbb Z}
\newcommand{\E}{\mathbb {E}}
\newcommand{\Pbb}[1]{\P\left( #1\right)}
\newcommand{\Ebb}[1]{\E\left[ #1\right]}

\newcommand{\lb}{\left (}
\newcommand{\rb}{\right )}
\newcommand{\lbb}{\left [}
\newcommand{\rbb}{\right ]}
\newcommand{\labs}{\left |}
\newcommand{\rabs}{\right |}
\newcommand{\lbrb}[1]{\lb #1 \rb}
\newcommand{\lbbrbb}[1]{\lbb#1\rbb}
\newcommand{\labsrabs}[1]{\labs#1\rabs}
\newcommand{\lbbrb}[1]{\lbb#1\rb}
\newcommand{\lbrbb}[1]{\lb#1\rbb}
\newcommand{\langlerangle}[1]{\langle#1\rangle}
\newcommand{\lbcurly}{\{}
\newcommand{\rbcurly}{\}}
\newcommand{\curly}[1]{\lbcurly#1\rbcurly}
\newcommand{\abs}{\labsrabs}
\newcommand{\ind}[1]{1_{#1}}
\newcommand{\indep}{\perp \!\!\! \perp}
\renewcommand{\d}{\, \mathrm{d}}

\newtheorem{theorem}{Theorem}
\newtheorem{problemCnt}{Lemma}
\theoremstyle{definition}
\newtheorem{problem}[problemCnt]{Problem}
\newtheorem{problembook}[problemCnt]{\faBook \,  Problem}
\newtheorem{problempencil}[problemCnt]{\faPencil* \,  Problem}
\newtheorem{problemstar}[problemCnt]{\faStar\,  Problem}
\newtheorem{problemstarpencil}[problemCnt]{\faStar \, \faPencil* \,   Problem}
\newtheorem{probleminteresting}[problemCnt]{\faSmile[regular] \,   Problem}
\newtheorem{solution}{Solution}

\newcommand{\fwcell}[1]{\makebox[\arraystretch\normalbaselineskip]{$#1$}}
\newenvironment{enumeratei}{\begin{enumerate}[(i)]}{\end{enumerate}}
\newenvironment{enumeratea}{\begin{enumerate}[(a)]}{\end{enumerate}}
\newenvironment{enumeraten}{\begin{enumerate}[\upshape 1.]}{\end{enumerate}}

% Compatibility helpers needed by a few files in the combined archive.
\AtBeginDocument{\noindent{\footnotesize English translation prepared with ChatGPT (OpenAI, GPT-6), 26 October 2026, for translation of the Bulgarian text into English and LaTeX formatting.
The Bulgarian originals remain the reference; translation may contain errors.}\par\smallskip}
\newcommand{\MissingGraphic}[1]{%
  \begin{center}\fbox{\parbox{0.82\linewidth}{\centering Missing source image: \texttt{#1}}}\end{center}}



\setlength{\emergencystretch}{3em}
\allowdisplaybreaks
\begin{document}
\noindent Software Engineering \hfill Test 1, 01.12.2024: solutions
\par\smallskip\hrule\smallskip\noindent
\textit{English transcription of the handwritten solutions. The source's order, calculations and unfinished evaluations are retained.}

\section*{Problem 1}
\begin{enumerate}
\item $\P(A)\P(B)=0.12$, and
\begin{align*}
\P(A\cap B)&=\P(A)+\P(B)-\P(A\cup B)\\
&=\P(A)+\P(B)-(1-\P(\overline{A\cup B}))\\
&=0.7-0.58=0.12.
\end{align*}
Thus $\P(A)\P(B)=\P(A\cap B)$: yes, they are independent.
\item
\[
\P(\text{three correct})=\binom53(\tfrac14)^3(\tfrac34)^2\approx8.79\%,
\]
\[
\P(\text{questions 1,3,5 correct}\mid\text{three correct})
=\frac1{\binom53}=\frac1{10}.
\]
\item Under independence, all the covariances are zero. In general, use the formula (or definition):
\begin{align*}
\operatorname{Cov}(X+Y,Z)
&=\E[(X+Y)Z]-\E[X+Y]\E Z\\
&=\E[XZ]-\E X\E Z+\E[YZ]-\E Y\E Z\\
&=\operatorname{Cov}(X,Z)+\operatorname{Cov}(Y,Z),
\end{align*}
by linearity.
\item Let $X_i=1$ if toss $i$ is heads and zero otherwise; the $X_i$ are iid $\operatorname{Ber}(1/2)$. We seek
\[
\operatorname{Cov}(X_1+\dots+X_7,(1-X_4)+\dots+(1-X_{10})).
\]
By part 3 this is $4\operatorname{Cov}(X_i,-X_i)=-4D X_i=-1$, since $\operatorname{Cov}(X_i,X_j)=0$ for $i\neq j$.
Writing the two sums as $Y,Z$, the correlation is
\[
\frac{-1}{\sqrt{DY\,DZ}}=\frac{-1}{7\cdot\frac14}=-\frac47.
\]
\item Ending after exactly three rounds means no ace among cards 1--20 and at least one ace among cards 21--30. Thus
\begin{align*}
\P(\text{exactly three rounds})
&=\P(\text{no ace in cards 1--20})-\P(\text{no ace in cards 1--30})\\
&=\frac{\binom{48}{20}}{\binom{52}{20}}-\frac{\binom{48}{30}}{\binom{52}{30}}\\
&=\frac{48!}{20!28!}\frac{20!32!}{52!}
-\frac{48!}{30!18!}\frac{30!22!}{52!}\\
&=\frac{29\cdot30\cdot31\cdot32}{49\cdot50\cdot51\cdot52}
-\frac{19\cdot20\cdot21\cdot22}{49\cdot50\cdot51\cdot52}
\approx10.58\%.
\end{align*}
\item Let $X_1,\dots,X_6$ be the face labels, iid uniform on $\{1,\dots,6\}$, and let $Y$ be the roll. Then
\begin{align*}
\P(\text{six})&=\E[\one_{\{\text{six}\}}]\\
&=\E\!\left[\E\!\left[\frac16\one_{\{X_1=6\}}+\dots+
\frac16\one_{\{X_6=6\}}\ \middle|\ X_1,\dots,X_6\right]\right]\\
&=\E[\one_{\{X_i=6\}}]=\frac16.
\end{align*}
Alternatively, condition on the number of sixes, which ranges from 0 to 6.
The expected number of rolls until the product is divisible by 6 is $\infty$ (or one may say it does not exist).
The die could, for example, have only ones; then the product can never be divisible by 6.
\end{enumerate}

\section*{Problem 3}
\begin{enumerate}
\item Coupon collector: see the exercises.
\[
\text{Answer: }10(1+\tfrac12+\dots+\tfrac1{10})\approx29.29.
\]
\item $X_i\sim\operatorname{Ge}(1/10)$ and $Y=\max\{X_1,\dots,X_{10}\}$.
\item
\[
\max\{x_1,\dots,x_n\}=
\sum_i x_i-\sum_{i<j}\min\{x_i,x_j\}
+\sum_{i<j<k}\min\{x_i,x_j,x_k\}-\dots
+(-1)^{n+1}\min\{x_1,\dots,x_n\}.
\]
\item $\operatorname{Ge}(p_1+p_2-p_1p_2)$. See the exercises.
\item For four species with probabilities $1/10,2/10,3/10,4/10$,
\[
\min\{X_i,X_j\}\sim\operatorname{Ge}(p_i+p_j).
\]
This does not follow from part 4, because the $X_i,X_j$ are not independent; it follows by another argument (what is it?).
Hence
\begin{align*}
\E[\max\{X_1,\dots,X_4\}]
&=\E[\sum_iX_i]-\E[\sum_{i<j}\min\{X_i,X_j\}]+\dots\\
&=\left[10+\frac{10}{2}+\frac{10}{3}+\frac{10}{4}\right]\\
&\quad-\left[\frac{10}{3}+\frac{10}{4}+\frac{10}{5}+
\frac{10}{5}+\frac{10}{6}+\frac{10}{7}\right]\\
&\quad+\left[\frac{10}{6}+\frac{10}{7}+\frac{10}{8}+\frac{10}{9}\right]-1=\dots
\end{align*}
\end{enumerate}

\section*{Problem 2}
\begin{enumerate}
\item Let $A_i$ mean that player $i$ receives BGN 0. Then
\begin{align*}
\P(A_1\cup A_2\cup A_3)
&=\P(\text{at least one receives 0})\\
&=\sum_i\P(A_i)-\sum_{i<j}\P(A_i\cap A_j)+\P(A_1\cap A_2\cap A_3)\\
&=3\P(A_1)-3\P(A_1\cap A_2)+\P(A_1\cap A_2\cap A_3)\\
&=3\sum_{\ell=0}^\infty(\tfrac12)^{\ell+1}
\left[(\tfrac23)^\ell-(\tfrac13)^\ell\right]+\frac12\\
&=\frac32\left[\frac1{1-\frac13}-\frac1{1-\frac16}\right]+\frac12\\
&=\frac32\left[\frac32-\frac65\right]+\frac12=\frac{19}{20}.
\end{align*}
The answer is $1-\P(\bigcup_iA_i)=1/20=5\%$.
\item Let $X$ be the winnings of a fixed player. For $k\geq0$,
\[
\P(X=k)=\sum_{\ell=0}^\infty
(\tfrac12)^{\ell+1}\binom\ell k(\tfrac13)^k(\tfrac23)^{\ell-k}.
\]
Here $(1/2)^{\ell+1}$ is the probability that the casino pays BGN $\ell$, and the remaining factor is the probability that $k$ of these payments go to the fixed player.
Put $\ell=k+m$:
\begin{align*}
\P(X=k)&=\sum_{m=0}^\infty(\tfrac12)^{k+m+1}
\binom{k+m}{k}(\tfrac13)^k(\tfrac23)^m\\
&=(\tfrac16)^k\frac12\sum_{m=0}^\infty(\tfrac13)^m\binom{k+m}{k}\\
&=(\tfrac16)^k\frac12\frac1{(1-\frac13)^{k+1}}
=\frac12(\tfrac16)^k(\tfrac32)^{k+1}
=\frac34(\tfrac14)^k.
\end{align*}
The sum uses the formula supplied in the test. The handwritten conclusion is
\[
X\sim\operatorname{Ge}(\tfrac14)
\quad\Longrightarrow\quad
DX=\frac{3/4}{(1/4)^2}=12.
\]
\end{enumerate}
\end{document}
